\section{Beyond the Condition Number}
\label{sec:beyond-kappa}

The condition number is a worst-case spectral parameter: it charges for a
small eigenvalue even when the right-hand side barely sees its eigenspace.
Dalzell, Li, and Su (DLS) recently gave two quantum linear-system solvers that
replace this worst-case dependence by right-hand-side-dependent truncation and
filtering parameters \cite{dalzell2026-faster-quantum-linear-system-solver};
their filtering solver improves an earlier construction of
\citet{li2025-new-condition-number-application-multivariate-polynomial}.
The classical idea of effective conditioning is older
\cite{chan1988-effectively-well-conditioned-linear-systems}, but the DLS
parameters had not been evaluated on IPM Newton systems.  We do so here.

Sections~\ref{sec:conditioning}--\ref{sec:spectra} showed that a degenerate LP
produces two increasingly separated spectral clusters.  The question in this
section is finer: when does the Newton right-hand side allow a solver to exploit
that separation?  Exact block identities give a complete coupling criterion.
They also expose a truncation transition, while a basis-invariant spectral
dimension gives more general solver-selection laws.  Finally, when the active
range is supplied, a directly assembled equilibration removes the central-path
divergence while preserving the Newton state exactly.

Throughout the first part, $H\succ0$ is presented through a block encoding of
$H/\alpha_H$.  For a nonzero right-hand side $f$, set $x=H^{-1}f$ and
define
\begin{equation}
 \xi(H,f)=\alpha_H\frac{\norm{H^{-1}f}}{\norm f},
 \qquad
 \rho(H,f)=\alpha_H
 \frac{\norm{H^{-2}f}}{\norm{H^{-1}f}}.
 \label{eq:dls-two-parameters}
\end{equation}
The second quantity is the ideal DLS filtering parameter: it is
$\norm{(H/\alpha_H)^{-1}x}/\norm x$, not the first-inverse
quantity $\xi$.  We avoid $\nu$, which \cref{sec:prelim-conic} reserves for
the barrier parameter.  More generally, for a nonsingular, possibly
non-Hermitian matrix $A$ encoded as $A/\alpha_A$, filtering uses
\begin{equation}
 \rho(A,f)=\alpha_A\frac{\norm{A^{-\dagger}x}}{\norm x},
 \qquad x=A^{-1}f.
 \label{eq:dls-general-filtering-parameter}
\end{equation}
The DLS truncation parameter instead asks how far toward zero
one must invert in order to retain all but the prescribed normalized solution
mass.

\subsection{Exact-center coupling law}
\label{sec:beyond-kappa-coupling}

Consider a standard-form LP with full-row-rank $A$, and suppose its exact
central path converges to a strictly complementary solution.  Let
\[
 B=\{i:x_i^*>0\},\qquad N=\{i:s_i^*>0\},\qquad
 U=\operatorname{range}(A_B),\qquad K=U^\perp,
\]
and assume $0<\dim U<m$.  This is the mixed, primal-degenerate case; if
$U=\R^m$, the normalized normal matrix has no asymptotically slow cluster.
At an exact center, complementarity gives
\begin{equation}
 H_\mu=H_{\rm NE}(\mu):=AXS^{-1}A^\top
 =\mu^{-1}C_\mu+\mu D_\mu,
 \quad
 C_\mu=A_BX_B^2A_B^\top,
 \quad
 D_\mu=A_NS_N^{-2}A_N^\top.
 \label{eq:coupling-normal-split}
\end{equation}
Here $C_\mu\to C_0$, $D_\mu\to D_0$, and, in $U\oplus K$ blocks,
\[
 D_\mu=
 \begin{pmatrix}D_\mu^U&F_\mu^\top\\F_\mu&E_\mu\end{pmatrix}.
\]
The restrictions $C=C_0|_U$ and $E=E_0|_K$ are positive definite.  Indeed,
$C$ is a positive-weight Gram matrix with range $U$, and if $z\in K$ obeys
$A_N^\top z=0$, then also $A_B^\top z=0$; full row rank of $A$ forces
$z=0$, proving positivity of $E$ on $K$.  At the optimal limit,
$b=Ax^*=A_Bx_B^*$, so the exact-centering right-hand side lies in $U$.  A
short step to
$\sigma_\mu\mu$, $\sigma_\mu\ne1$, has
$H_\mu\Delta y=(1-\sigma_\mu)b$.  The scalar is immaterial to normalized
QLS parameters, so the next theorem uses $b\ne0$. This nonzero condition
also follows in the stated LP setting. Indeed $x^*\ne0$; if $b=0$,
optimality gives $c^Tx^*=0$, while a strictly feasible dual slack $s$
would give $s^Tx^*=c^Tx^*>0$, a contradiction.

\begin{theorem}[Exact-center coupling criterion]
\label{thm:dls-coupling}
Define
\begin{align}
 R_\mu&=D_\mu^U-F_\mu^\top E_\mu^{-1}F_\mu,\nonumber\\
 p_\mu&=(C_\mu|_U+\mu^2R_\mu)^{-1}b,
 &g_\mu&=E_\mu^{-1}F_\mu p_\mu,\label{eq:coupling-pg}\\
 q_\mu&=(C_\mu|_U+\mu^2R_\mu)^{-1}
 (p_\mu+F_\mu^\top E_\mu^{-1}g_\mu).\nonumber
\end{align}
Then the following identities are exact:
\begin{align}
 H_\mu^{-1}b
 &=\mu\begin{pmatrix}p_\mu\\-g_\mu\end{pmatrix},
 \label{eq:coupling-first-inverse}\\
 H_\mu^{-2}b
 &=\begin{pmatrix}
       \mu^2q_\mu\\
       -E_\mu^{-1}g_\mu-\mu^2E_\mu^{-1}F_\mu q_\mu
    \end{pmatrix}.
 \label{eq:coupling-second-inverse}
\end{align}
For a norm-tight encoding,
$\alpha_H=\Theta(\norm{H_\mu})=\Theta(\mu^{-1})$,
\begin{equation}
 \xi(H_\mu,b)=\Theta(1),\qquad
 \rho(H_\mu,b)=\Theta\!\left(1+\frac{\norm{g_\mu}}{\mu^2}\right).
 \label{eq:coupling-filtering-law}
\end{equation}
In particular, with
\begin{equation}
 g=E^{-1}FC^{-1}b,
 \label{eq:limiting-coupling}
\end{equation}
filtering improves asymptotically on the full condition-number scale if and
only if $g=0$.  If $g\ne0$, then
$\rho(H_\mu,b)=\Theta(\mu^{-2})=\Theta(\kappa(H_\mu))$; if $g=0$, then
$\rho=O(\mu^{-1})=O(\sqrt{\kappa(H_\mu)})$, with bounded cost precisely when
$g_\mu=O(\mu^2)$.
\end{theorem}

\begin{proof}
Write a vector in $U\oplus K$ coordinates.  Multiplying
$H_\mu(u,v)^\top=b$ by $\mu$, the lower block equation gives
$v=-E_\mu^{-1}F_\mu u$.  Substitution in the upper equation gives
$(C_\mu|_U+\mu^2R_\mu)u=\mu b$, proving
\eqref{eq:coupling-first-inverse}.

Apply the same elimination to
$H_\mu(u,v)^\top=H_\mu^{-1}b$.  Its lower equation is
$F_\mu u+E_\mu v=-g_\mu$, so
$v=-E_\mu^{-1}g_\mu-E_\mu^{-1}F_\mu u$.  The upper equation then becomes
\[
 (C_\mu|_U+\mu^2R_\mu)u
 =\mu^2(p_\mu+F_\mu^\top E_\mu^{-1}g_\mu),
\]
which proves \eqref{eq:coupling-second-inverse}.

Continuity and positive definiteness give
$p_\mu\to p=C^{-1}b\ne0$.  Moreover,
$q_\mu\to C^{-1}(p+F^\top E^{-1}g)$, and the vector in parentheses is
nonzero because its inner product with $p$ is
$\norm p^2+\norm{E^{-1}Fp}^2>0$.  Equations
\eqref{eq:coupling-first-inverse}--\eqref{eq:coupling-second-inverse}
therefore imply
\[
 \norm{H_\mu^{-1}b}=\Theta(\mu),\qquad
 \norm{H_\mu^{-2}b}=\Theta(\norm{g_\mu}+\mu^2).
\]
For the lower bound, write the second inverse as $(a_\mu,v_\mu)$,
where $a_\mu=\mu^2q_\mu$. Its lower block equation gives
$g_\mu=-E_\mu v_\mu-F_\mu a_\mu$. Boundedness of $E_\mu,F_\mu$
controls $\norm{g_\mu}$ by $\norm{(a_\mu,v_\mu)}$, and
$q_\mu\to q\ne0$ controls $\mu^2$ by the same norm. This also covers
cancellation between the two terms in the lower block.
Substitution in \eqref{eq:dls-two-parameters} proves
\eqref{eq:coupling-filtering-law}.  Finally $g_\mu\to g$.  Thus $g\ne0$
gives the full $\mu^{-2}$ scale.  Under strict complementarity the LP
central path, and hence $C_\mu,D_\mu,F_\mu,E_\mu$ and $g_\mu$, is analytic
at $\mu=0$.  Halick\'a proves this boundary analyticity for LP and the
corresponding result for SDP
\cite{halicka1999-analytical-properties-central-path-linear-programming,
halicka2002-analyticity-of-the-central-path}. If $g=0$, differentiability
alone gives $g_\mu=O(\mu)$, hence $\rho=O(\mu^{-1})$. More precisely,
let $g'_0$ denote the derivative of the coupling at zero. If $g'_0\ne0$,
then $\norm{g_\mu}=\Theta(\mu)$ and $\rho=\Theta(\mu^{-1})$.
If $g'_0=0$, the analytic Taylor remainder gives $g_\mu=O(\mu^2)$
and bounded $\rho$. Thus boundedness is equivalent to $g_0=g'_0=0$
for this analytic family. The norm law itself also proves boundedness
exactly when $g_\mu=O(\mu^2)$.
\end{proof}

Equivalently, the obstruction is
$\Pi_KD_0(C|_U)^{-1}b\ne0$.  The exceptional limiting right-hand-side space
is $C(\ker F)$, and the obstruction vanishes for every $b\in U$ exactly when
$F=0$, meaning that $D_0$ reduces $U\oplus K$.  Analyticity sharpens the
continuous-rate analysis: after $g=0$, the only divergent intermediate rate
is a first-order term $g_\mu=\Theta(\mu)$, which gives
$\rho=\Theta(\mu^{-1})$; every higher integer order gives bounded $\rho$.
For a non-tight encoding, the exact scale carried by the proof is
\begin{equation}
 \rho(H_\mu,b)=\Theta\!\left(
 \alpha_H\left(\mu+\frac{\norm{g_\mu}}\mu\right)\right).
 \label{eq:coupling-nontight}
\end{equation}

\begin{proposition}[Sparse sharpness witnesses]
\label{prop:coupling-witnesses}
There are two row- and column-two LPs with the same
$\kappa(H_\mu)\sim(2\mu^2)^{-1}$ and $\xi=\Theta(1)$, but respectively
$\rho=\Theta(\mu^{-2})$ and $\rho=1$.
\end{proposition}

\begin{proof}
For the coupled instance, take
\begin{equation}
 \min\ x_2+x_3
 \quad\text{subject to}\quad
 \begin{pmatrix}1&1&0\\0&1&-1\end{pmatrix}x
 =\begin{pmatrix}1\\0\end{pmatrix},\qquad x\ge0.
 \label{eq:coupled-small-lp}
\end{equation}
Its primal optimum and optimal partition are unique.  The optimal duals have
$y_1=0$, $y_2\in[-1,1]$; their strictly complementary members have
$y_2\in(-1,1)$.  The displayed slack
$s^*=(0,1,1)$ corresponds to $y^*=0$, the dual analytic center and hence the
central-path limit, while $x^*=e_1$.  The central point is
$x=(1-t,t,t)$, where
\[
 t=\frac{2+3\mu-\sqrt{4-4\mu+9\mu^2}}4=\mu+O(\mu^2).
\]
With $\theta_1=(1-t)^2/\mu$ and $\theta=t^2/\mu$,
\[
 H_\mu=
 \begin{pmatrix}\theta_1+\theta&\theta\\\theta&2\theta\end{pmatrix},
 \qquad \kappa(H_\mu)\sim\frac1{2\mu^2}.
\]
Here $C=1$, $E=2$, $F=1$, and the centering right-hand side is
$b=(1,0)^\top$, so $p=1$ and $g=1/2$.  Direct inversion gives
\[
 H_\mu^{-1}b=\frac{(2,-1)^\top}{2\theta_1+\theta},
 \qquad H_\mu^{-2}b\longrightarrow(0,-1/4)^\top.
\]
For $\alpha_H=\norm{H_\mu}$,
$\xi\to\sqrt5/2$ and $\mu^2\rho\to1/(2\sqrt5)$.  The normalized solution
is exactly $(2,-1)/\sqrt5$. Since the slow eigenspace tends to the
second coordinate axis, its slow-cluster amplitude tends to $1/\sqrt5$.

For the decoupled instance, replace the constraint matrix by
\[
 \begin{pmatrix}1&0&0\\0&1&-1\end{pmatrix}
\]
and keep the same objective and right-hand side.  It has the same unique
primal optimum, optimal partition, dual-optimal interval, central-path dual
limit, and it also satisfies the proposition's row/column-two sparsity bound,
but, for $0<\mu\le1/\sqrt2$ and $\alpha_H=\norm{H_\mu}$,
\[
 H_\mu=\diag(\mu^{-1},2\mu),\qquad
 \kappa(H_\mu)=\frac1{2\mu^2},\qquad \xi=\rho=1.
\]
Thus neither $\kappa$ nor the first-inverse norm distinguishes the two
instances; active-to-inactive coupling does.
\end{proof}

\paragraph{Formal verification and scope.}
The Lean development in \path{formal/QipmFormal/Coupling/} checks the LP
normal-matrix split and centering right-hand side, block elimination and
the actual inverse identities, noncancellation, and the exceptional
right-hand-side space. It derives the filtering law from continuous block
data and checks the differentiable and analytic rate consequences.
The sparse witnesses are checked from their primal and dual equations,
with Euclidean norms and the displayed normalization qualifications.
The claim map and reproduction command are in \path{formal/COUPLING.md}.
The audit permits only \texttt{propext}, \texttt{Classical.choice}, and
\texttt{Quot.sound}. The general existence and boundary analyticity of
the LP central path remain mathematical inputs; the development verifies
the block regularity consequences rather than formalizing those external
theorems. It does not verify the quantum solver, construct its block
encoding, or establish the later truncation and spectral-dimension claims.

\subsection{Truncation at the slow-shell boundary}
\label{sec:beyond-kappa-truncation}

Let $P_s(\mu)$ project onto the eigenvalues of $H_\mu$ of order
$\Theta(\mu)$, which become $\Theta(\mu^2)$ after norm normalization.
For a general nonzero $f$, define its slow right-hand-side amplitude and
slow solution mass by
\[
 q_s=\frac{\norm{P_sf}}{\norm f},\qquad
 p_s=\frac{\norm{P_sH_\mu^{-1}f}^2}{\norm{H_\mu^{-1}f}^2}.
\]

\begin{proposition}[Two-cluster truncation transition]
\label{prop:dls-truncation-transition}
If the fast-cluster component of $f$ is bounded below, then
\begin{equation}
 p_s=\Theta\!\left(\frac{q_s^2}{\mu^4+q_s^2}\right).
 \label{eq:slow-solution-mass}
\end{equation}
Consequently $q_s=\Theta(\mu^2)$ is the critical boundary.  There the slow
solution mass is constant and fixed-accuracy DLS truncation must retain the
full $\Theta(\mu^{-2})$ normalized inverse scale for every target error below
that mass.  If $q_s=o(\mu^2)$, truncation can eventually discard the slow
cluster at fixed accuracy; filtering need not yet be bounded.  Under a tight
encoding, $q_s=O(\mu^4)$ is sufficient for bounded ideal filtering cost.
At exact centering, generic coupling gives $q_s=\Theta(\mu^2)$.
\end{proposition}

\begin{proof}
On the unnormalized fast and slow clusters, $H_\mu^{-1}$ has scales
$\Theta(\mu)$ and $\Theta(\mu^{-1})$, respectively.  The lower bound on the
fast part of $f$ therefore gives fast solution norm $\Theta(\mu\norm f)$,
while the slow solution norm is $\Theta(\mu^{-1}q_s\norm f)$.  Dividing the
squared slow norm by the sum proves \eqref{eq:slow-solution-mass}.

For the exact-centering vector $b\in U$, standard separated-subspace
perturbation, or direct substitution in the block eigenvector equation,
gives in the leading $K$ component
\begin{equation}
 P_s(\mu)b=-\mu^2D_{0,KU}C^{-1}b+o(\mu^2).
 \label{eq:exact-center-slow-rhs}
\end{equation}
It is therefore $\Theta(\mu^2)$ when the limiting coupling is nonzero.
Equation \eqref{eq:slow-solution-mass} proves the truncation claims.  One
additional inverse weights the slow part by another $\mu^{-2}$ relative to
the fast part, so its contribution is bounded only at
$q_s=O(\mu^4)$, which is the filtering statement.
\end{proof}

The comparison with \cref{prop:newton-rhs-alignment,cor:top-window-solve}
requires naming the two formulations.  This section uses the $m$-dimensional
normal matrix
\[
 H_{\rm NE}(\mu)=AXS^{-1}A^\top,
 \qquad
 \operatorname{spec}(H_{\rm NE}(\mu))
 \subset\Theta(\mu^{-1})\cup\Theta(\mu),
\]
whereas \cref{sec:spectra-rhs} uses the $(n-m)$-dimensional reduced Hessian
\[
 H_{\rm red}(\mu)=Z^\top X^{-2}Z,
 \qquad
 \operatorname{spec}(H_{\rm red}(\mu))
 \subset\Theta(\mu^{-2})\cup\Theta(1).
\]
When the corresponding clusters are nonempty, both have fast normalized
scale $\Theta(1)$ and slow normalized scale $\Theta(\mu^2)$, but they act on
different variables and right-hand sides.  \Cref{sec:spectra} proves an
$O(\mu^2)$ slow right-hand-side fraction for $H_{\rm red}(\mu)$.  For
$H_{\rm NE}(\mu)$, the present coupling calculation
instead gives $q_s=\Theta(\mu^2)$ only when the limiting coupling is nonzero;
the decoupled witness has $P_sb=0$.

The residual-versus-state distinction can be proved directly for the normal
equation.  Under nonzero limiting coupling, let
$z=H_{\rm NE}(\mu)^{-1}b$ and discard
$z_s=P_sz$.  Because $P_s$ is a spectral projector,
\begin{equation}
 \frac{\norm{H_{\rm NE}(\mu)(z-z_s)-b}}{\norm b}
 =\frac{\norm{P_sb}}{\norm b}=q_s=\Theta(\mu^2),
 \label{eq:normal-window-residual}
\end{equation}
while \eqref{eq:slow-solution-mass} gives
$\norm{z_s}^2/\norm z^2=\Theta(1)$ under the same nonzero-coupling
hypothesis.  A state-preparation algorithm below that constant solution mass
must retain the slow cluster, but a fixed relative-residual contract may
discard it.  Thus the two formulations agree on their normalized spectral
separation, while both the formulation and the output contract matter; this
is the precise sense in which the result is consistent with
\cref{sec:spectra}.

\subsection{Effective spectral dimension}
\label{sec:effective-spectral-dimension}

Two clusters are not necessary for a sharp solver-selection law.  Normalize
$\norm H=1$ and $\norm b=1$, let
\begin{equation}
 \mathsf m_b(S)=\inner b{\mathbf1_S(H)b},
 \label{eq:spectral-measure}
\end{equation}
and let $\lambda_*$ be the smallest eigenvalue visible to $b$, with
$\kappa_*=\lambda_*^{-1}$.  We say that $(H,b)$ has effective spectral
dimension $D$ if constants $c,C>0$, independent of $\lambda_*$, satisfy
\begin{equation}
 c t^{D/2}\leq F_b(t):=\mathsf m_b([0,t])\leq C t^{D/2}
 \qquad(\lambda_*\leq t\leq1).
 \label{eq:effective-dimension-definition}
\end{equation}
This is invariant under simultaneous orthogonal changes
$(H,b)\mapsto(VHV^\top,Vb)$; it is not coordinate sparsity in disguise.
Unlike an unanchored per-shell condition, \eqref{eq:effective-dimension-definition}
controls the mass at the smallest visible eigenvalue and permits fixed gaps
between the first few eigenvalues.

\begin{theorem}[Spectral-dimension phase laws]
\label{thm:spectral-dimension-laws}
Assume the norm-tight unit-normalization convention: after $\norm H=1$, the
block-encoding normalization is $\alpha_H=\Theta(1)$, and truncation cutoffs
are measured in this same normalization.  Operationally,
\begin{equation}
 \kappa_{\rm eff}(\epsilon_{\rm ls})
 :=\inf\left\{K\in[1,\kappa_*]:
 \frac{\inner b{H^{-2}\mathbf1_{[0,1/K)}(H)b}}{\inner b{H^{-2}b}}
 \leq\epsilon_{\rm ls}^2\right\}.
 \label{eq:dls-effective-cutoff-definition}
\end{equation}
For an encoding of $H/\alpha_H$, the solver's encoded cutoff scale is
$\alpha_H\kappa_{\rm eff}$; the standing tight convention identifies these
up to constants.
Under \eqref{eq:effective-dimension-definition}, this quantity obeys
\begin{equation}
 \kappa_{\rm eff}(\epsilon_{\rm ls})=
 \begin{cases}
  \Theta_{\epsilon_{\rm ls}}(\kappa_*),&D<4,\ 0<\epsilon_{\rm ls}<1,\\
  \kappa_*^{\,1-\Theta(\epsilon_{\rm ls}^2)},&D=4,\\
  \Theta\!\left(\min\{\kappa_*,
                  \epsilon_{\rm ls}^{-4/(D-4)}\}\right),&D>4.
 \end{cases}
 \label{eq:dls-truncation-dimension-law}
\end{equation}
For $D<4$, the hidden constant may depend on the fixed error and on the
cumulative-law constants.  A sublinear power of $\kappa_*$ is possible only
when the allowed error tends to one as $\kappa_*\to\infty$.  At $D=4$, a
definite leading cumulative coefficient gives exponent
$1-\epsilon_{\rm ls}^2+o(1)$; under only
\eqref{eq:effective-dimension-definition},
$1-\Theta(\epsilon_{\rm ls}^2)$ is sharp, and the result saturates at
$\Theta(\kappa_*)$ when
$\epsilon_{\rm ls}^2\log\kappa_*=O(1)$.  The ideal filtering
parameter is
\begin{equation}
 \rho=
 \begin{cases}
  \Theta(\kappa_*),&D<4,\\
  \Theta(\kappa_*/\sqrt{\log\kappa_*}),&D=4,\\
  \Theta(\kappa_*^{(8-D)/4}),&4<D<8,\\
  \Theta(\sqrt{\log\kappa_*}),&D=8,\\
  \Theta(1),&D>8.
 \end{cases}
 \label{eq:dls-filter-dimension-law}
\end{equation}
Thus fixed-accuracy truncation becomes independent of mesh conditioning above
$D=4$, whereas filtering does so only above $D=8$.  For a non-tight encoding,
both the effective cutoff scale and every filtering branch acquire a factor
$\alpha_H$; in particular the $D>8$ filtering cost is
$\Theta(\alpha_H)$ rather than $\Theta(1)$.
\end{theorem}

\begin{proof}
For $p>0$, put $Z_p=\inner b{H^{-p}b}$ and write $s=D/2$.  Abel summation,
equivalently Stieltjes integration by parts, gives the exact layer-cake
identity
\begin{equation}
 Z_p=1+p\int_{\lambda_*}^1 t^{-p-1}F_b(t)\,dt.
 \label{eq:spectral-abel-identity}
\end{equation}
Using $F_b(t)=\Theta(t^s)$ in this integral gives
\begin{equation}
 Z_p=
 \begin{cases}
  \Theta(\kappa_*^{p-D/2}),&D<2p,\\
  \Theta(\log\kappa_*),&D=2p,\\
  \Theta(1),&D>2p.
 \end{cases}
 \label{eq:spectral-moment-law}
\end{equation}
The squared norm of $H^{-1}b$ is $Z_2$.  If a cutoff discards eigenvalues
strictly below $u=1/K$, as in \eqref{eq:dls-effective-cutoff-definition}, its
discarded squared solution norm is
\begin{equation}
 T_2(u)=u^{-2}F_b(u^-)
       +2\int_{\lambda_*}^{u}t^{-3}F_b(t)\,dt,
 \label{eq:spectral-partial-abel}
\end{equation}
where $F_b(u^-)=\mathsf m_b([0,u))$ is the left limit.
For $D<4$, fix a sufficiently small constant $u_0>0$.  Dyadically summing
the retained moment on $[u,1]$ and comparing it with $Z_2$ gives
\begin{equation}
 \frac{\inner b{H^{-2}\mathbf1_{[u,1]}(H)b}}{Z_2}
 =\Theta\!\left(
   \left(\frac{\lambda_*}{u}\right)^{2-D/2}\right)
 \qquad(\lambda_*\leq u\leq u_0).
 \label{eq:subcritical-retained-mass}
\end{equation}
Indeed, the upper bound follows by geometric summation, while a shell
$[u,\gamma_{\rm sh}u]$ with a fixed $\gamma_{\rm sh}>1$ chosen from the two
cumulative-law constants has mass $\Theta(u^{D/2})$ and gives the matching
lower bound.  Requiring the
retained fraction to be at least $1-\epsilon_{\rm ls}^2$ therefore gives
\begin{equation}
 u=\Theta\!\left(
 \lambda_*(1-\epsilon_{\rm ls}^2)^{-2/(4-D)}\right)
 =\Theta_{\epsilon_{\rm ls}}(\lambda_*),
 \label{eq:subcritical-cutoff}
\end{equation}
for every fixed $\epsilon_{\rm ls}<1$, proving the first line.  A cutoff
$u=\lambda_*^{1-\beta}$ for fixed $\beta>0$ has retained fraction tending to
zero, so it is
possible only when $\epsilon_{\rm ls}\to1$.  For $D=4$, both
\eqref{eq:spectral-abel-identity} and \eqref{eq:spectral-partial-abel} are
logarithmic.  Discarding an $\epsilon_{\rm ls}^2$ fraction yields
$K=\kappa_*^{1-\Theta(\epsilon_{\rm ls}^2)}$, with the stated refinement when the
leading cumulative coefficient exists.  For $D>4$, the discarded tail is
$\Theta(K^{-(D-4)/2})$ relative to the constant total.  Setting it to
$\epsilon_{\rm ls}^2$ and capping at $\kappa_*$ proves
\eqref{eq:dls-truncation-dimension-law}.

For filtering, $x=H^{-1}b$ and hence
\[
 \rho^2=\alpha_H^2\frac{\norm{H^{-1}x}^2}{\norm x^2}
 =\alpha_H^2\frac{Z_4}{Z_2}=\Theta\!\left(\frac{Z_4}{Z_2}\right).
\]
Applying \eqref{eq:spectral-moment-law} with $p=4$ in the numerator and
$p=2$ in the denominator gives all five cases in
\eqref{eq:dls-filter-dimension-law}.
\end{proof}

Local conventions in the remaining witnesses are as follows:
$e_0,e_{e_1},e_f$ are coordinate basis vectors, whereas $e_k$ in the
marked-diagonal construction is the $k$-dimensional all-ones vector; the
symbols $\delta_{\rm c},\delta_Q,\delta_{\mathcal G}$ denote, respectively,
the corrector perturbation, projector error, and expander contrast.

\begin{proposition}[Constant-degree torus witnesses]
\label{prop:torus-dimension-witnesses}
Let $L$ be even and $T_L^d$ the periodic $d$-dimensional $L$-grid.  With
its combinatorial Laplacian $\Delta_{T_L^d}$, set
\begin{equation}
 H_{L,d}=\frac{L^{-2}I+\Delta_{T_L^d}}{L^{-2}+4d}.
 \label{eq:massive-torus}
\end{equation}
Then $H_{L,d}$ is $(2d+1)$-sparse, has norm one, and has
$\kappa=1+4dL^2=\Theta(L^2)$.  A localized source $e_0$ realizes $D=d$;
the dipole $(e_0-e_{e_1})/\sqrt2$ realizes $D=d+2$.  For the dipole branch,
the resulting phases are:
\begin{center}
\begin{tabular}{@{}rrrr@{}}
\toprule
physical $d$ & effective $D$ & $\kappa_{\rm eff}$ & $\rho$ \\
\midrule
$3$ & $5$ & $\Theta(\min\{\kappa,\epsilon_{\rm ls}^{-4}\})$
            & $\Theta(\kappa^{3/4})$ \\
$4$ & $6$ & $\Theta(\min\{\kappa,\epsilon_{\rm ls}^{-2}\})$
            & $\Theta(\kappa^{1/2})$ \\
$5$ & $7$ & $\Theta(\min\{\kappa,\epsilon_{\rm ls}^{-4/3}\})$
            & $\Theta(\kappa^{1/4})$ \\
$6$ & $8$ & $\Theta(\min\{\kappa,\epsilon_{\rm ls}^{-1}\})$
            & $\Theta(\sqrt{\log\kappa})$ \\
$\geq7$ & $>8$ & mesh-independent at fixed $\epsilon_{\rm ls}$
                   & mesh-independent \\
\bottomrule
\end{tabular}
\end{center}
\end{proposition}

\begin{proof}
Fourier modes diagonalize the torus Laplacian, with
\[
 \lambda_k=\frac{L^{-2}+4\sum_{j=1}^d\sin^2(\pi k_j/L)}{L^{-2}+4d}.
\]
For the localized source, every squared Fourier overlap is $L^{-d}$.
Lattice-volume bounds, first for
$\lambda_k\asymp L^{-2}(1+\norm k^2)$ at low frequency and then by the
same product-counting bounds on the remaining fixed spectral range, give
\[
 \mathsf m_{e_0}([0,t])=\Theta(t^{d/2})
 \qquad(\lambda_*\leq t\leq1).
\]
Thus the cumulative hypothesis, including its bottom anchoring, holds with
$D=d$.  Equivalently, Abel summation gives
\[
 \sum_k\lambda_k^{-p}=
 \begin{cases}
  \Theta(L^{2p}),&d<2p,\\
  \Theta(L^{2p}\log L),&d=2p,\\
  \Theta(L^d),&d>2p,
 \end{cases}
\]
The dipole Fourier symbol is $1-e^{2\pi i k_1/L}$.  It vanishes on modes
with $k_1=0$, so no pointwise comparison with $\lambda_k$ is valid.  Instead
sum over permutation-symmetric spectral balls.  Symmetry gives
\[
 \sum_{\lambda_k\leq t}\sin^2(\pi k_1/L)
 =\frac1d\sum_{\lambda_k\leq t}\sum_{j=1}^d
       \sin^2(\pi k_j/L).
\]
Together with the lattice-volume estimate, this aggregated identity yields
\[
 \mathsf m_{(e_0-e_{e_1})/\sqrt2}([0,t])
 =\Theta(t^{(d+2)/2})
\]
from the smallest visible eigenvalue through $t=1$.  At the bottom, the
first visible modes have squared overlap $\Theta(L^{-(d+2)})$, exactly
$\Theta(\lambda_*^{(d+2)/2})$; the constant mode has zero overlap and
$\lambda_*=\Theta(L^{-2})$.  Hence the repaired cumulative definition is
satisfied with $D=d+2$.  The table follows by substituting $D=d+2$ in
\cref{thm:spectral-dimension-laws}.
\end{proof}

The two source types occur as actual QIPM correctors, not merely abstract
linear systems.

\begin{proposition}[Positive full-step feasible corrector]
\label{prop:torus-feasible-corrector}
Let $B_T$ be an oriented incidence matrix of $T_L^d$ and set
\begin{equation}
 \mathcal A=[B_T,I],\qquad
 W=\diag(I_E,L^{-2}I_V),\qquad
 \mathcal A W\mathcal A^\top=\Delta_{T_L^d}+L^{-2}I.
 \label{eq:torus-ipm-factor}
\end{equation}
For $d\geq2$, choose perpendicular torus edges $f,g$ meeting at a vertex,
fix $\tau>0$ and $0<\delta_{\rm c}\leq1/4$, and retain $x_i/s_i=w_i$, where
$w_i$
is the corresponding diagonal entry of $W$.  Set
\begin{equation}
 p_i=x_is_i=
 \begin{cases}
  \tau(1+\delta_{\rm c}),&i=f,\\
  \tau(1-\delta_{\rm c}),&i=g,\\
  \tau,&\text{otherwise},
 \end{cases}
 \quad
 x_i=\sqrt{p_iw_i},\quad s_i=\sqrt{p_i/w_i}.
 \label{eq:torus-perturbed-products}
\end{equation}
With $b_{\rm LP}=\mathcal Ax$, $c=s$, and $y=0$, this is primal and dual
feasible, has average complementarity $\tau$, and has centrality
$\norm{XSe-\tau e}/\tau=\sqrt2\delta_{\rm c}$.  Its full $\sigma=1$ feasible
centering step is positive, has a normal-equation right-hand side of
effective dimension $D=d+2$, preserves average complementarity, and satisfies
\begin{equation}
 \frac{\norm{X^+S^+e-\tau e}}\tau
 \leq\frac{\delta_{\rm c}^2}{1-\delta_{\rm c}^2}.
 \label{eq:torus-quadratic-centering}
\end{equation}
\end{proposition}

\begin{proof}
The definitions make feasibility and the centrality calculation immediate.
The feasible corrector equations are
\[
 \mathcal A\Delta x=0,
 \quad \mathcal A^\top\Delta y+\Delta s=0,
 \quad S\Delta x+X\Delta s=\tau e-XSe.
\]
Eliminating $\Delta x,\Delta s$ gives
\begin{equation}
 (\mathcal AW\mathcal A^\top)\Delta y
 =-\mathcal AS^{-1}(\tau e-XSe).
 \label{eq:torus-corrector-normal}
\end{equation}
After normalization, its right-hand side is proportional to
\[
 \frac{\mathcal A_f}{\sqrt{1+\delta_{\rm c}}}
 -\frac{\mathcal A_g}{\sqrt{1-\delta_{\rm c}}}.
\]
Its low-frequency Fourier symbol vanishes linearly and has nonzero gradient,
so its cumulative spectral weight is $\Theta(t^{(d+2)/2})$.  Explicitly, for
any nonzero
linear form $\ell$, lattice-ball comparison gives
$\sum_{\norm k\leq r}|\ell(k)|^2=\Theta(r^{d+2})$; applying this to the
linear part of the symbol gives the claimed cumulative spectral weight.

It remains to prove full-step safety.  Put
$u=X^{-1}\Delta x$, $v=S^{-1}\Delta s$, and
$P=\diag(p)$.  The Newton
equations imply
\begin{equation}
 P(u+v)=\tau e-p,\qquad
 u^\top Pv=\Delta x^\top\Delta s=0.
 \label{eq:torus-relative-step}
\end{equation}
Thus $P^{1/2}u$ and $P^{1/2}v$ are orthogonal and sum to
$P^{-1/2}(\tau e-p)$, whose squared norm is
$2\tau\delta_{\rm c}^2/(1-\delta_{\rm c}^2)$.  Since
$p_i\geq\tau(1-\delta_{\rm c})$, for either $z=u,v$,
\[
 |z_i|^2\leq
 \frac{2\delta_{\rm c}^2}
 {(1-\delta_{\rm c}^2)(1-\delta_{\rm c})}<1
 \qquad(\delta_{\rm c}\leq1/4).
\]
Hence $x+\Delta x>0$ and $s+\Delta s>0$.  Moreover
\[
 p_i^+=p_i(1+u_i)(1+v_i)=\tau+p_i u_iv_i.
\]
Equation \eqref{eq:torus-relative-step} gives
$\sum_i p_iu_iv_i=0$, so the average remains $\tau$.  Finally,
\[
 \norm{P(u\circ v)}
 \leq\norm{P^{1/2}u}\norm{P^{1/2}v}
 \leq\frac12\norm{P^{-1/2}(\tau e-p)}^2
 =\frac{\tau\delta_{\rm c}^2}{1-\delta_{\rm c}^2},
\]
which proves \eqref{eq:torus-quadratic-centering}.
\end{proof}

The localized branch is also QIPM-native, but through an infeasible
corrector.  At complementarity $\tau$, set
$x_i=\sqrt{\tau w_i}$, $s_i=\sqrt{\tau/w_i}$ and
$b_{\rm LP}=\mathcal Ax$.  Give the exactly centered
($x_is_i=\tau>0$), primal-feasible
point a one-coordinate dual residual $r_d=\eta e_f$.  A $\sigma=1$
infeasible feasibility corrector has
\begin{equation}
 (\mathcal AW\mathcal A^\top)\Delta y=\eta\mathcal AWe_f.
 \label{eq:torus-infeasible-corrector}
\end{equation}
A ground variable at vertex zero gives normalized right-hand side $e_0$ and
$D=d$; a torus edge gives a dipole and $D=d+2$.

These witnesses are solver-selection examples, not end-to-end speedups.  The
torus already admits FFT and multigrid solvers.  The constructions describe
a single corrector rather than an IPM trajectory, and do not charge block
encoding, norm estimation, state or observable readout, or classical iterate
recovery.

\subsection{Stable active-subspace equilibration}
\label{sec:stable-active-equilibration}

The negative coupling law has a positive structural counterpart.  Suppose
the active range is available through its orthogonal projector $Q$, and put
$R=I-Q$.  The next theorem directly assembles a left-equilibrated operator;
this direct assembly, rather than multiplication of normalized encodings, is
the source of the gain.

\begin{theorem}[Stable split equilibration]
\label{thm:stable-split-equilibration}
Suppose
\[
 H_\mu=\mu^{-1}C_\mu+\mu D_\mu,\qquad 0<\mu\leq1,
\]
where, for constants independent of $\mu$,
\[
 C_\mu=QC_\mu Q,\quad
 c_-Q\preceq C_\mu\preceq c_+Q,\quad
 D_\mu\succeq0,\quad \norm{D_\mu}\leq d_+,\quad
 RD_\mu R\succeq d_-R.
\]
Let $P_\mu=\mu Q+\mu^{-1}R$ and assemble
\begin{equation}
 M_\mu=C_\mu+RD_\mu+\mu^2QD_\mu
 =C_\mu+D_\mu-(1-\mu^2)QD_\mu.
 \label{eq:stable-split-operator}
\end{equation}
Then $M_\mu=P_\mu H_\mu$ exactly and
$\norm{M_\mu}+\norm{M_\mu^{-1}}=O(1)$.  For every
$f\in U=\operatorname{range}Q$,
\begin{equation}
 M_\mu^{-1}f=\mu^{-1}H_\mu^{-1}f,
 \label{eq:stable-state-preservation}
\end{equation}
so the normalized Newton state is preserved exactly, without recovery or
postselection.  If $C_\mu,D_\mu,Q$ have block-encoding normalizations
$\alpha_C,\alpha_D,1$, direct LCU assembly has normalization
\begin{equation}
 \Gamma_\mu\leq\alpha_C+(2-\mu^2)\alpha_D.
 \label{eq:stable-lcu-normalization}
\end{equation}
Thus bounded component normalizations give a uniformly conditioned,
uniformly normalized conventional QLS solve through the Hermitian dilation
below, and the ideal DLS filtering parameter is $O(\Gamma_\mu)$ as well.
\end{theorem}

The left-equilibrated $M_\mu$ need not be Hermitian.  To use the Hermitian
QLS contract of \cref{def:qls}, encode its standard dilation
\begin{equation}
 \mathcal H(M_\mu)=
 \begin{pmatrix}0&M_\mu\\M_\mu^\dagger&0\end{pmatrix}.
 \label{eq:stable-hermitian-dilation}
\end{equation}
The dilation has the same block-encoding normalization $\Gamma_\mu$ and
the same singular condition number as $M_\mu$.  Moreover,
$\mathcal H(M_\mu)^{-1}(f,0)^\top=(0,M_\mu^{-1}f)^\top$, so selecting the
second register preserves the solution normalization and requires no
postselection beyond the fixed register label.

\begin{proof}
Since $RC_\mu=0$,
\[
 P_\mu H_\mu=C_\mu+\mu^2QD_\mu+RD_\mu=M_\mu.
\]
In $U\oplus U^\perp$ blocks,
\[
 M_\mu=
 \begin{pmatrix}
 C_\mu+\mu^2D_{UU}&\mu^2D_{UK}\\
 D_{KU}&D_{KK}
 \end{pmatrix}.
\]
The Schur complement of $D_{KK}$ is
\[
 G_\mu=C_\mu+\mu^2
 (D_{UU}-D_{UK}D_{KK}^{-1}D_{KU})\succeq c_-I_U,
\]
because the parenthesized Schur complement is positive semidefinite.
Also $D_{KK}\succeq d_-I_{U^\perp}$.  The block inverse formula therefore
bounds every block of $M_\mu^{-1}$ using only $c_-,d_-,d_+$, while
$\norm{M_\mu}\leq c_++2d_+$.

For state preservation, let $x_\mu=H_\mu^{-1}f$.  Because $f\in U$,
$P_\mu f=\mu f$, and hence
$M_\mu x_\mu=P_\mu H_\mu x_\mu=\mu f$.  This proves
\eqref{eq:stable-state-preservation}.  Finally, applying LCU to the second
form in \eqref{eq:stable-split-operator} gives coefficient weight
$\alpha_C+\alpha_D+(1-\mu^2)\alpha_D$, proving
\eqref{eq:stable-lcu-normalization}.  Uniform conditioning bounds
$\Gamma_\mu\norm{M_\mu^{-1}}$ and its restriction to any solution state,
which proves the two solver claims.
\end{proof}

The split also supplies a projector-error firewall.  If
$\norm{\widehat Q-Q}\leq\delta_Q$ and one assembles
\[
 \widehat M_\mu=C_\mu+D_\mu-(1-\mu^2)\widehat QD_\mu,
\]
then $\norm{\widehat M_\mu-M_\mu}\leq\delta_Qd_+$.  For a sufficiently small
constant $\delta_Q$, the Neumann perturbation bound and
\cref{thm:stable-split-equilibration} give
\begin{equation}
 \frac{\norm{\widehat M_\mu^{-1}f-M_\mu^{-1}f}}
      {\norm{M_\mu^{-1}f}}=O(\delta_Q),
 \label{eq:projector-firewall}
\end{equation}
uniformly in $\mu$.  Thus $\delta_Q=O(\epsilon_{\rm ls})$ suffices.  In
contrast, forming $\widehat P_\mu H_\mu$ from the raw normal matrix loses
this stability.  For $H_\mu=\diag(\mu^{-1},\mu)$ and a projector rotated by
angle $\theta$, its $(2,1)$ entry is
$(1-\mu^{-2})\sin\theta\cos\theta$; uniform normalization requires
$\sin\theta=O(\mu^2)$.  Given a block encoding of $A_B/\alpha_B$, QSVT can
construct $Q$ to the needed accuracy in
\begin{equation}
 O\!\left(\frac{\alpha_B}{\sigma_{\min}^+(A_B)}
 \log(1/\epsilon_{\rm ls})\right)
 \label{eq:active-projector-cost}
\end{equation}
queries, charging active-basis conditioning but no polynomial power of
$\mu$.

\begin{proposition}[Separately normalized composition cannot help]
\label{prop:no-composed-preconditioning}
Let $H/\alpha_H$ and an invertible $P/\alpha_P$ be separately encoded, with
$\alpha_H\geq\norm H$ and $\alpha_P\geq\norm P$, and compose their encodings.
For $x=H^{-1}f$, the DLS parameter of the resulting system
$PHx=Pf$ satisfies
\begin{equation}
 \rho_{\rm prod}
 =\alpha_H\alpha_P
  \frac{\norm{P^{-\dagger}H^{-1}x}}{\norm x}
 \geq \alpha_H\frac{\norm{H^{-1}x}}{\norm x}
 =\rho_H.
 \label{eq:no-composition-law}
\end{equation}
\end{proposition}

\begin{proof}
Equation \eqref{eq:no-composition-law} uses the general non-Hermitian
parameter \eqref{eq:dls-general-filtering-parameter}, because
$(PH)^{-\dagger}=P^{-\dagger}H^{-1}$.  Since
$\alpha_P\geq\norm P$, every
vector $y$ satisfies
$\alpha_P\norm{P^{-\dagger}y}\geq\norm y$.  Apply this with
$y=H^{-1}x$.  The displayed expression is precisely the filtering
parameter for normalization $\alpha_H\alpha_P$, proving the result.
\end{proof}

This normalization obstruction is known in quantum preconditioning:
Lapworth and S\"underhauf show explicitly how multiplying separately
normalized block encodings can erase a classical gain, while Nie, Lai, and
An show how algebraic regrouping can restore it for Pauli structure
\cite{leigh2025-preconditioned-block-encodings-quantum-linear,hantao2026-pauli-structured-preconditioning-quantum-linear}.
Here separate composition has $\alpha_H,\alpha_P=\Theta(\mu^{-1})$, whereas the
direct split has $\Gamma_\mu=O(1)$, a $\Theta(\mu^{-2})$ normalization
compression.

\begin{proposition}[Constant-sparse expander occurrence]
\label{prop:stable-expander-witness}
There is a family of constant-sparse standard-form LPs for which the
unpreconditioned exact-centering solve has $\rho=\Theta(\mu^{-2})$, while
\eqref{eq:stable-split-operator} has $O(1)$ singular conditioning and
$O(1)$ block-encoding normalization and prepares exactly the same normalized
Newton state independently of $\mu$.  The uniform bounds hold on
$0<\mu\leq\mu_0$ with $\mu_0$ and all constants depending only on the degree
$d$, the Laplacian gap, and $\delta_{\mathcal G}$, not on $m$.
\end{proposition}

\begin{proof}
Let $\mathcal G=(\mathcal L,\mathcal R,\mathcal E)$ be a connected balanced
$d$-regular bipartite expander on $m$ vertices, with constant $d$ and a
constant Laplacian gap.  Orient each edge from $\mathcal L$ to $\mathcal R$,
let $B_{\mathcal G}$ be its signed incidence matrix, and put
\[
 r=\mathbf1_{\mathcal R},\qquad \zeta=2r-e,\qquad
 u_0=m^{-1/2}e.
\]
Here $r$ is the right-part indicator and $e$ is the all-ones vector of the
dimension required by context; these symbols are local to this witness and
are distinct from the coupling variables above.  Then
$B_{\mathcal G}^\top r=e$ and
$B_{\mathcal G}e=L_{\mathcal G}r=d\zeta$.  Fix
$0<\delta_{\mathcal G}<1$, set
$\omega_i=1+\delta_{\mathcal G}\zeta_i$ and
$a_i=\sqrt{2/\omega_i}$, and consider
\begin{equation}
 \begin{aligned}
 \min_{x,u,v\geq0}\quad&\sum_{i=1}^m a_i(u_i+v_i)\\
 \text{subject to}\quad&B_{\mathcal G}x+u-v=B_{\mathcal G}e.
 \end{aligned}
 \label{eq:expander-lp}
\end{equation}
Its constraint matrix $[B_{\mathcal G},I,-I]$ has row and column sparsity
$O(d)$.

The point $x=e,u=v=te$ is strictly primal feasible for $t>0$.  Taking
$y=-\varepsilon r$ for sufficiently small $\varepsilon>0$
gives strictly positive dual slacks.  Nonnegativity of the objective shows
that
\[
 x^*=e,\quad u^*=v^*=0,\quad y^*=0,
 \qquad s_x^*=0,\quad s_u^*=s_v^*=a
\]
is optimal and strictly complementary.  Its active face is bounded: if
$h\geq0$ and $B_{\mathcal G}h=0$, any left-vertex row forces all incident
components of $h$ to vanish.  Also $e$ is the primal face's analytic center
because $B_{\mathcal G}^\top r=e$.

For completeness, the dual limit is also unique even though the dual optimal
face is not a singleton.  Dual feasibility and optimality imply
$B_{\mathcal G}^\top y\leq0$ and
$e^\top B_{\mathcal G}^\top y=0$, hence
$B_{\mathcal G}^\top y=0$.  Connectedness gives $y=ce$, and the remaining
dual inequalities give
\[
 \{y:\ y\text{ is dual optimal}\}
 =\{ce:\ |c|\leq\min_i a_i\}.
\]
On its relative interior the dual analytic-center objective is
$\sum_i\log(a_i^2-c^2)$, which is strictly maximized at $c=0$.  Thus
$y(\mu)\to0$, and the displayed primal--dual point is the central-path
limit.

Along the path,
\begin{align}
 H_\mu&=\mu^{-1}C_\mu+\mu D_\mu,\nonumber\\
 C_\mu&=B_{\mathcal G}\diag(x_{\mathcal E}(\mu)^2)
          B_{\mathcal G}^\top\longrightarrow L_{\mathcal G},\nonumber\\
 D_\mu&=\diag\!\left(
 \frac1{(a_i-y_i(\mu))^2}+\frac1{(a_i+y_i(\mu))^2}
 \right)\longrightarrow\diag(\omega_i).
 \label{eq:expander-normal-split}
\end{align}
Here $U=\operatorname{range}B_{\mathcal G}=e^\perp$ and
$Q=I-\ket{u_0}\bra{u_0}$.  The constants are uniform in $m$.  Write
$s_x=-B_{\mathcal G}^\top y$, $s_u=a-y$, $s_v=a+y$; the central equations
are $x_es_{x,e}=\mu$, $u_is_{u,i}=\mu$, $v_is_{v,i}=\mu$, and
$B_{\mathcal G}x+u-v=B_{\mathcal G}e$.  Seek a solution constant on edges
and on each side: $x_e=\bar x$, $y_i=y_{\mathcal L}$ on $\mathcal L$,
$y_i=y_{\mathcal R}$ on $\mathcal R$, with $\bar x=1+\mu^2w$,
$y_{\mathcal L}=\mu z_{\mathcal L}$,
$y_{\mathcal R}=\mu z_{\mathcal R}$.  Since
$\omega=1\mp\delta_{\mathcal G}$ on $\mathcal L$ and $\mathcal R$, the
system reduces to three scalar equations independent of $m$:
\[
 dw+\frac{2z_{\mathcal R}}
 {2/(1+\delta_{\mathcal G})-\mu^2z_{\mathcal R}^2}=0,
 \qquad
 -dw+\frac{2z_{\mathcal L}}
 {2/(1-\delta_{\mathcal G})-\mu^2z_{\mathcal L}^2}=0,
 \qquad
 z_{\mathcal L}-z_{\mathcal R}=\frac1{1+\mu^2w}.
\]
At $\mu=0$ they have the solution
$w=(1-\delta_{\mathcal G}^2)/(2d)$,
$z_{\mathcal L}=(1+\delta_{\mathcal G})/2$,
$z_{\mathcal R}=-(1-\delta_{\mathcal G})/2$, with Jacobian determinant
$-2d\neq0$.  The implicit function theorem gives a branch on
$0<\mu\le\mu_0(d,\delta_{\mathcal G})$ with $\bar x>0$,
$y_{\mathcal L}-y_{\mathcal R}=\mu/\bar x>0$, and $|y_i|<a_i$; the ansatz
therefore satisfies the full central-path system, and uniqueness of the
central point identifies it as the central path.  Hence
$C_\mu=\bar x^2L_{\mathcal G}$ exactly, $D_\mu$ takes two side-constant
values, and the expander gap together with
$\omega_i\in[1-\delta_{\mathcal G},1+\delta_{\mathcal G}]$ gives the
hypotheses of \cref{thm:stable-split-equilibration} with constants depending
only on $d$, the Laplacian gap, and $\delta_{\mathcal G}$.

For the exact-centering right-hand side $f=b=B_{\mathcal G}e$,
$L_{\mathcal G}(Qr)=b$ and $Qr=\zeta/2$.  Its limiting active-to-null
response is nonzero:
\begin{equation}
 u_0^\top\diag(\omega)Qr=\frac{\delta_{\mathcal G}\sqrt m}{2}
 =\delta_{\mathcal G}\norm{Qr}.
 \label{eq:expander-coupling}
\end{equation}
The coupling theorem therefore gives $\rho(H_\mu,f)=\Theta(\mu^{-2})$.
In contrast, direct assembly gives
\begin{equation}
 M_\mu=C_\mu+\mu^2D_\mu
 +(1-\mu^2)\ket{u_0}\bra{u_0}D_\mu.
 \label{eq:expander-direct-split}
\end{equation}
It has uniform singular conditioning and $O(d)$ normalization: $C_\mu$ is
a bounded-weight degree-$d$ Laplacian, $D_\mu$ is diagonal, $u_0$ is the
uniform state, and the last term is rank one.  Since $b\in U$,
$M_\mu^{-1}b=\mu^{-1}H_\mu^{-1}b$ at every $\mu$, completing the witness.
\end{proof}

The theorem assumes a supplied active partition, coherent current weights,
and direct access to $C_\mu,D_\mu,Q$.  It also charges
$\sigma_{\min}^+(A_B)$, right-hand-side preparation, norm estimation, and
the requested output.  \Cref{sec:positive-modules} gives a partition-free
construction on
the strict-complementarity tail and thereby removes the supplied-partition
assumption; that construction is not needed here.  Classical application of
\eqref{eq:expander-direct-split} is already linear-time per Krylov iteration.
The quantum statement is therefore meaningful only for a state or a few
observables; tomography or a classical iterate restores dimension-dependent
cost.

\subsection{An exact-center marked-diagonal specialization}
\label{sec:marked-diagonal-specialization}

The sparse search lower bound of Orsucci and Dunjko
\cite[Proposition~7]{orsucci2021-on-solving-classes-of-positive} can also be
realized at an exact LP center.  This gives a useful QIPM specialization,
but its marked-search core is prior art.

Choose an integer $T\geq2$, set the local dimension
$N_{\rm md}=T^8$ and $\mu=T^{-2}=N_{\rm md}^{-1/4}$, and hide
$j\in[N_{\rm md}]$.  The subscript distinguishes this dimension from the
optimal-partition set $N$ above.  Let $e_k$ denote the $k$-dimensional
all-ones vector in this construction.  Row $i$ of
$A_j\in\R^{N_{\rm md}\times2N_{\rm md}}$ has disjoint support
$(2i-1,2i)$ and values
\begin{equation}
 (\mu/2,\mu/2)\quad(i=j),\qquad
 (1,\mu-1)\quad(i\ne j).
 \label{eq:marked-diagonal-A}
\end{equation}
Because $\mu$, $\sqrt\mu$, and $\mu^{3/2}$ are rational, all entries have
magnitude at most one and $O(\log N_{\rm md})$-bit rational descriptions;
row sparsity is two and column sparsity one.  Set
\[
 b=\mu^{3/2}e_{N_{\rm md}},\qquad
 c=\sqrt\mu e_{2N_{\rm md}},\qquad
 x^0=s^0=\sqrt\mu e_{2N_{\rm md}},\qquad y^0=0.
\]
Then $A_jx^0=b$, $(A_j)^\top y^0+s^0=c$, and
$X^0S^0e=\mu e$, so this is an exact $\mu$-center.  Since
$X^0(S^0)^{-1}=I$,
\begin{equation}
 H_j=A_jA_j^\top
 =\diag(h,\ldots,\mu^2/2,\ldots,h),
 \qquad h=1+(1-\mu)^2,
 \label{eq:marked-diagonal-H}
\end{equation}
and the exact short-step right-hand side is
$f=(1-\sigma)\mu^{3/2}e_{N_{\rm md}}$.  Its normalized preparation is uniform
and independent of $j$.  The squared solution weight on the marked
coordinate is at least $4/5$, while
\[
 \begin{aligned}
 \kappa(H_j)=\Theta(\mu^{-2}),\qquad
 \xi(H_j,f)=\Theta(1),\qquad
 \rho(H_j,f)&=\Theta(\mu^{-2}),\\
 \kappa_{\rm eff}(\epsilon_{\rm ls})
 &=\frac{2h}{\mu^2}=\Theta(\mu^{-2})
 \quad\text{for every fixed }\epsilon_{\rm ls}^2<4/5.
 \end{aligned}
\]
After norm-tight normalization the only eigenvalues are $1$ and
$\mu^2/(2h)$, so $\kappa_{\rm eff}$ drops to $1$ only once
$\epsilon_{\rm ls}^2$ reaches the marked weight, which is at least $4/5$.
Consequently, under the stated sparse-value oracle, or the canonical
diagonal block encoding reducible to it, preparing the conventional exact
normalized multiplier state to trace error $1/10$ (so that measuring the
marked coordinate succeeds with probability at least $4/5-1/10=7/10$)
requires
$\Omega(\sqrt{N_{\rm md}})=\Omega(\mu^{-2})=\Omega(\kappa)$ queries.

Taking $1-\sigma=\mu^2/4=\Theta(N_{\rm md}^{-1/2})$ gives a positive feasible
short step whose endpoint remains in a fixed two-norm central neighborhood.  This
makes the occurrence native to an IPM, subject to three essential caveats.
First, it is not representation invariant: row equilibration turns
\eqref{eq:marked-diagonal-H} into the identity and makes the scaled
multiplier state easy (the primal/slack direction itself is invariant).
Second, simply omitting the marked component leaves relative normal residual
only $\Theta(\mu^2)$, so ordinary residual-based inexact IPMs can bypass the
state-preparation contract.  Third, the hard scale occurs at one point of
each fixed instance's path; it does not persist as that fixed instance's
barrier parameter tends to zero.  The lower bound is therefore for the
original exact normalized multiplier-state contract, not for every QIPM or
every equivalent representation.  This distinction also accords with the
representation and finite-precision cautions in classical IPM analysis
\cite{wright1998-ill-conditioning-computational-error-interior}.
This $\Omega(\kappa)$ statement is a full-query-model bound and survives
factor access.  It is consistent with
\cref{thm:quantum-cluster-obstruction}\textup{(iii)}, which is a
polynomial-metric statement and becomes an end-to-end factor-access solve
only under its explicit overlap hypothesis.

\subsection{Novelty and scope}

The application of the 2026 DLS solvers to IPM Newton systems, the exact
coupling law, the $D=4$ versus $D=8$ selection rule and its QIPM correctors,
and the state-preserving stable split appear to be new.  This is deliberately
not an unconditional priority claim.  Effective conditioning, active-set
preconditioning, Green-function moment thresholds, and the generic
marked-diagonal search bound are prior art.  None of the results in this
section proves an end-to-end quantum LP speedup or a new universal QLS lower
bound.  Access construction, active-set discovery where needed, precision,
iterate maintenance, and output costs remain part of any such claim.
